\documentclass[10pt,a4paper]{article}
\usepackage{amsmath,amssymb,amsthm,graphicx}
\usepackage{amsthm,amsmath}
\usepackage{natbib,setspace}
\usepackage{color,comment}

\newtheorem{theorem}{Theorem}
\newtheorem{corollary}{Corollary}
\newtheorem{lemma}{Lemma}
\newtheorem{assumption}{Assumption}
\newtheorem{example}{Example}
\newtheorem{fact}{Fact}

\usepackage{amsmath,amssymb,amsfonts}

\usepackage{natbib}
 \bibpunct[, ]{(}{)}{,}{a}{}{,}%
 \def\bibfont{\small}%
 \def\bibsep{\smallskipamount}%
 \def\bibhang{24pt}%
 \def\newblock{\ }%
 \def\BIBand{and}%

\usepackage{rotating}
\usepackage{fancyvrb}
\usepackage{mathrsfs}
\newcommand{\ve}{\mathscr{V}}
\title{Wasserstein convergence rate for stationary finite-state vector Markov additive functionals, allowing periodicity}
\author{R. Srikant}
\date{arxiv:2401.15719v5; CC BY4.0}
\begin{document}
\maketitle
\noindent\textit{Editorial note.} The fixed source is arXiv2401.15719v5 dated2026-08-29. Source123 explicitly states that this version corrects errors in the Math of OR published version; do not substitute a prior version or reuse its proof silently. The officialv5 source archive names its main filearxiv-v6.tex, which is an internal filename rather than the arXiv version number.\par\medskip
\noindent\textit{Editorial note.} Only the finite-state stationary theorem946-962 is counted. Periodicity is allowed; the separate arbitrary-initial-distribution corollary and its aperiodicity assumption are not included in this claim.\par\medskip
\noindent\textit{Editorial note.} The already public abstract martingale theorem199-281 and its complete proof are included as supporting solution obligations. No new count is assigned to that theorem or to the Hessian-variation lemma separately.\par\medskip
\noindent\textit{Editorial note.} The source uses the same letterh when passing through the Lipschitz rescaling bijection305-311 and taking the supremum at738. The final theorem is a supremum statement; the review must preserve that rescaling convention rather than identify the unscaled and normalized individual test functions.\par\medskip
\noindent\textit{Editorial scope.} One stationary finite-state vector MarkovWassersteinCLT and its optimized log-squared rate, with all local core retained. The general-state extension and TD-learning theorem have not been fully prepared by this card.\par\medskip
\noindent\textit{Allowed standard prerequisites (editorial).}\begin{itemize}
\item Finite irreducible Markov-chain Poisson equation. For a finite irreducible transition matrixP and a centered functiong withpi(g)=0, (I-P)V=g has a bounded solution, also for a periodic chain; apply entrywise to centered matrix functions.. Produce reward and covariance potentials. No quantitative Markov CLT is supplied.
\item Standard Gaussian Stein solution regularity. For a Lipschitz test function, the normal Stein solution has bounded Hessian and beta-Holder Hessian with constants of order(1-beta)\textasciicircum{}(-1), as explicitly stated in source.. Bound the two martingale Taylor errors. The local martingale bound, its signed covariance term and its control are proof obligations.
\item Gaussian integration by parts, Mehler formula and smoothing derivative estimates. ForP\_Ah(x)=E[h(x+A\textasciicircum{}(1/2)Z)] andh1-Lipschitz, ||nabla\textasciicircum{}j P\_Ah||<=C\_d||A\textasciicircum{}(-1/2)||\textasciicircum{}(j-1), j=2,3,4; use the Gaussian covariance interpolation identity.. Derive the source Hessian-weight representation and separately control shift/internal-covariance/external-covariance changes.
\item Martingale conditioning, finite-dimensional matrix inequalities and summation by parts. No independence between the random Stein Hessian and the Markov conditional covariance is assumed.
\end{itemize}\medskip
\subsection{Preliminaries}

In this subsection, we will describe the notation we use to establish the rate of convergence in the central limit theorem for vector-valued martingale differences. Let $\{m_k\}_{k\geq 1}$ be a $d$-dimensional martingale difference sequence with respect to a filtration $\{\mathcal{F}_k\}_{k\geq 0}$ i.e., $m_k$ is $\mathcal{F}_k$-adapted, $E(||m_k||)<\infty,$ where $||\cdot||$ is the standard Euclidean norm in $\Re^d$ and $$E(m_k|\mathcal{F}_{k-1})=0\qquad \forall k\geq 1.$$ 
\begin{assumption}\label{assume: martingale}
Fix $\beta\in(0,1)$. We assume that the martingale difference sequence satisfies the following properties:
\begin{enumerate}\label{assume}
\item $E\|m_k\|^{2+\beta}<\infty$ for every $k\geq 1$.
\item $\Sigma_k=E(m_km_k^T\mid\mathcal{F}_{k-1})$ exists for every $k\geq 1$.
\end{enumerate}
\hfill $\diamond$
\end{assumption}
Item (1) in the assumption is motivated by the fact that we want to extend the techniques in \citep{gallouet2018regularity} to martingales .

Define
$S_n=\sum_{k=1}^n m_k.$ We are interested in whether the following normalized sum 
\begin{equation}\label{eq: normalized sum}
W_n:=\frac{S_n}{\sqrt{n}}
\end{equation}
converges to $N(0,\Sigma_\infty)$, and if so, we are interested in the rate of convergence. We will use Stein's method to establish such a convergence and to estimate the rate of convergence \citep{stein1972bound,ross2011fundamentals,chatterjee2014short}. 

The Wasserstein distance between two $d$-dimensional random vectors $X$ and $Y$ is defined as
$$d_{\mathcal{W}}(X,Y)=\sup_{h\in Lip_1} E[h(X)-h(Y)],$$
where $Lip_L$, the class of Lipschitz functions from $\Re^d$ to $\Re$  with Lipschitz constant $L,$ is defined as
$$Lip_L=\{h: |h(x)-h(y)|\leq L||x-y||\}.$$
For later reference, we use the notation $||A||_{op}$  to denote the operator norm of a matrix $A,$ i.e., $$||A||_{op}=\sup_{x: ||x||=1}||Ax||.$$
The following results from \cite[Propositions 2.2 and 2.3]{gallouet2018regularity} and \cite[Theorem 3.1]{fang2019multivariate} will be useful to us and so we present them together as a lemma for future reference.
\begin{lemma}\label{lem: regularity}
Let $g\in Lip_L$ and $f$ be the solution to Stein's equation:
\begin{equation}\label{eq: Stein}
    g(u)-E(g(Z))=\Delta f(u)-u^T\nabla f(u),
\end{equation}
where $\Delta$ and $\nabla$ are the Laplacian and gradient operators, respectively, and $Z\sim N(0,I)$. Then, $f$ has the following regularity properties for any $\beta\in (0,1)$:
\begin{enumerate}
    \item $||\nabla^2f(x)-\nabla^2f(y)||_{op}\leq \tilde{C}_1(d,\beta)||x-y||^\beta L,$ where $\nabla^2$ is the Hessian operator.
    \item $|\Delta f(x)-\Delta f(y)|\leq \tilde{C_2}(d,\beta)||x-y||^\beta L$,
where  $$\tilde{C}_1=C_1(d)+\frac{2}{1-\beta}, \quad \tilde{C}_2=C_2(d)+\frac{2d}{1-\beta}, \quad 
C_1=2^{3/2}\frac{(1+2d)\Gamma((1+d)/2)}{d\Gamma(d/2)},\quad C_2=2\sqrt{\frac{2d}{\pi}}.$$
    \item  $$\left|\langle\nabla^2 f(x),u_1u_2^T\rangle_{HS}\right|\leq C L||u_1||\cdot||u_2||\,\forall x, u_1, u_2\in \Re^d,$$ where $C$ is a constant (independent of $d$ and $\beta$) and $\langle\cdot,\cdot\rangle_{HS}$ is the Hilbert-Schmidt inner product (also called the Frobenius inner product), i.e., the sum of the element-by-element product of two matrices.
\end{enumerate}
\hfill $\diamond$
\end{lemma}

\section{Markov Chain Central Limit Theorem}\label{sec: Markov}
Consider a time-homogeneous Markov chain $\{X_k\}_{k\geq 0}$ on a state space $\mathcal{S}$ with initial distribution $\eta$. Let $r(X_k)\in\Re^d$ denote a vector reward and let $\bar r=\pi r$ whenever the chain has invariant distribution $\pi$. We are interested in rates of convergence to normality of the centered additive functional
\[
    \frac{1}{\sqrt n}\sum_{k=1}^n\bigl(r(X_{k-1})-\bar r\bigr).
\]
As is common in the literature, we study this scaled sum by relating the reward process to a martingale via Poisson's equation \citep{douc2018markov,makowski2002poisson,glynn1996liapounov,glynn2023solution}.

\subsection{Finite State Space Markov Chains}
\label{subsec: finite_state_MC_corrected}

We now show that the martingale theorem can be used to obtain a
non-asymptotic CLT for finite-state Markov chains.  In Theorem~\ref{thm: martingale CLT}, one must control
the signed covariance-interaction term.  This can
be done by applying a second Poisson equation, now to the conditional
covariance function.

Throughout this subsection, assume that the Markov chain is stationary,
finite-state, and irreducible.  Periodicity is allowed.  Let
\[
    g(x):=r(x)-\bar r,
    \qquad
    \bar r:=\pi r.
\]
Let \(V\) solve the Poisson equation
\[
    V-PV=g.
\]
Then
\[
    g(X_{k-1})
    =
    V(X_{k-1})-V(X_k)+m_k,
\]
where
\[
    m_k:=V(X_k)-E[V(X_k)\mid X_{k-1}]
\]
is a martingale difference sequence.  Define
\[
    \Sigma_k:=E[m_km_k^\top\mid \mathcal F_{k-1}]
    =
    v(X_{k-1}),
\]
where
\[
    v(x):=E[m_1m_1^\top\mid X_0=x].
\]
Since the chain is stationary,
\[
    E(\Sigma_k)=\pi(v)=:\Sigma_\infty.
\]
Assume throughout this subsection that \(\Sigma_\infty\) is positive definite.

\begin{theorem}\label{thm: MC CLT}
Let
\[
U_n=\frac1{\sqrt n}\sum_{k=1}^n(r(X_{k-1})-\bar r).
\]
Assume that the Markov chain is stationary, finite-state, and irreducible,
and that $\Sigma_\infty\succ0$. Then, for every $\beta\in(0,1)$,
\[
d_{\mathcal W}(U_n,\Sigma_\infty^{1/2}Z)
\le \frac{C}{(1-\beta)^2}n^{-\beta/2},
\]
where $C$ does not depend on $n$ or $\beta$. Consequently,
\[
d_{\mathcal W}(U_n,\Sigma_\infty^{1/2}Z)
=O\left(\frac{(\log n)^2}{\sqrt n}\right).
\]
\end{theorem}

\begin{theorem}\label{thm: martingale CLT}
Let
\[
    V_k:=\sum_{j=k}^n E(\Sigma_j),
    \qquad k=1,\ldots,n,
\]
and assume that \(V_k\) is positive definite for each \(k=1,\ldots,n\).  Define
\[
    D_k:=V_k^{-1/2}E(\Sigma_k)V_k^{-1/2},
    \qquad k=1,\ldots,n .
\]
Let \(Z\sim N(0,I)\), and let \(W_n\) be defined as in
\((\ref{eq: normalized sum})\). Then, under
Assumption~\ref{assume: martingale} for the same fixed \(\beta\in(0,1)\),
\begin{align*}
   d_{\mathcal W}(W_n,n^{-1/2}V_1^{1/2}Z)
   &\le
   \frac{1}{\sqrt n}\sum_{k=1}^n
   \Biggl[
   \tilde C_1(d,\beta)\|V_k^{1/2}\|_{op}
   E\left(\|V_k^{-1/2}m_k\|^{2+\beta}\right)
   \\
   &\qquad\qquad
   +
   \tilde C_1(d,\beta)\operatorname{Tr}(D_k)\|V_k^{1/2}\|_{op}
   E\left(\|V_k^{-1/2}m_k\|^\beta\right)
   \Biggr]
   +\Delta_n,
\end{align*}
where
\begin{align*}
\Delta_n
:=
\sup_{h\in Lip_1}
\frac{1}{\sqrt n}
\left|
\sum_{k=1}^n
E\left[
\operatorname{Tr}\left(
A_{k,h}(\mathcal F_{k-1})
V_k^{-1/2}
\left(E(\Sigma_k)-\Sigma_k\right)
V_k^{-1/2}
\right)
\right]
\right|.
\end{align*}
Here
\[
A_{k,h}(\mathcal F_{k-1})
:=
E\left[
\nabla^2 f_{k,h}(\widetilde Z_k)
\mid
\mathcal F_{k-1}
\right],
\]
where \(f_{k,h}\) is the $\mathcal F_{k-1}$-measurable random solution of
the Stein equation defined in the proof below with
\[
    a=S_{k-1},
    \qquad
    Q_k=V_k^{1/2}.
\]
and
\[
    \widetilde Z_k
    :=
    V_k^{-1/2}T_{k+1}
    \sim
    N\left(0,V_k^{-1/2}V_{k+1}V_k^{-1/2}\right),
\]
with the convention \(V_{n+1}=0\). Moreover,
\[
    \|A_{k,h}(\mathcal F_{k-1})\|_{op}
    \le
    C\|V_k^{1/2}\|_{op}, \quad \text{a.s.},
\]
where \(C\) is the constant in statement (3) of Lemma~\ref{lem: regularity}.
Since \(0\preceq D_k\preceq I\), the same bound also implies the coarser
version obtained by replacing \(\operatorname{Tr}(D_k)\) by \(d\).  The sharper
form above is used in the Markov-chain applications.
\end{theorem}

\begin{proof}
Fix \(h\in Lip_1\). By the standard rescaling of Lipschitz test functions,
\[
d_{\mathcal W}(W_n,n^{-1/2}V_1^{1/2}Z)
=
\sup_{h\in Lip_1}
\frac{1}{\sqrt n}
E\left[h(S_n)-h(T_1)\right],
\]
where
\[
    T_k:=\sum_{j=k}^n Y_j,
\]
and \(Y_1,\ldots,Y_n\) are independent Gaussian random vectors, independent
of the martingale difference sequence, with
\[
    Y_j\sim N(0,E(\Sigma_j)).
\]
Thus
\[
    T_k\sim N(0,V_k).
\]
The rescaling follows because if \(g\in Lip_1\), then
\(h(x):=\sqrt n\,g(x/\sqrt n)\) also belongs to \(Lip_1\), and
\[
E\left[g(W_n)-g(n^{-1/2}V_1^{1/2}Z)\right]
=
\frac{1}{\sqrt n}
E\left[h(S_n)-h(T_1)\right].
\]

Using the Lindeberg decomposition,
\[
E\left[h(S_n)-h(T_1)\right]
=
\sum_{k=1}^n
E\left[
h(S_k+T_{k+1})-h(S_{k-1}+T_k)
\right],
\]
where \(T_{n+1}=S_0=0\).

We now bound a single summand. Fix \(k\). Conditional on
\(\mathcal F_{k-1}\), set
\[
    a=S_{k-1},
    \qquad
    Q_k=V_k^{1/2}.
\]
Define
\[
    \widetilde h_k(u):=h(a+Q_k u).
\]
Since \(h\in Lip_1\), we have
\[
    \widetilde h_k\in Lip_{L_k},
    \qquad
    L_k=\|Q_k\|_{op}=\|V_k^{1/2}\|_{op}.
\]
Let \(f_k=f_{k,h}\) be the solution of the Stein equation
\[
    \widetilde h_k(u)-E[\widetilde h_k(Z)]
    =
    \Delta f_k(u)-u^\top\nabla f_k(u).
\]
Since \(a=S_{k-1}\) is \(\mathcal F_{k-1}\)-measurable, the function
\(f_k\) is also \(\mathcal F_{k-1}\)-measurable as a random function.

By the Stein equation, with
\[
    \widetilde Z_k:=V_k^{-1/2}T_{k+1},
\]
we have
\begin{align*}
& E\left[
h(S_k+T_{k+1})-h(S_{k-1}+T_k)
\mid
\mathcal F_{k-1}
\right]
\\
&=
E\biggl[
\Delta f_k\left(V_k^{-1/2}m_k+\widetilde Z_k\right)
-
\left(V_k^{-1/2}m_k+\widetilde Z_k\right)^\top
\nabla f_k\left(V_k^{-1/2}m_k+\widetilde Z_k\right)
\mid
\mathcal F_{k-1}
\biggr].
\end{align*}
Moreover,
\[
    \widetilde Z_k
    \sim
    N(0,\Lambda_k),
    \qquad
    \Lambda_k:=V_k^{-1/2}V_{k+1}V_k^{-1/2}.
\]
Since
\[
    V_k=E(\Sigma_k)+V_{k+1},
\]
we have
\[
    I-\Lambda_k
    =
    V_k^{-1/2}E(\Sigma_k)V_k^{-1/2}.
\]

We rewrite the preceding conditional expectation as
\begin{align*}
& E\biggl[
\operatorname{Tr}\left(
\Lambda_k
\nabla^2 f_k\left(V_k^{-1/2}m_k+\widetilde Z_k\right)
\right)
-
\widetilde Z_k^\top
\nabla f_k\left(V_k^{-1/2}m_k+\widetilde Z_k\right)
\mid
\mathcal F_{k-1}
\biggr]
\\
&+
E\biggl[
\operatorname{Tr}\left(
(I-\Lambda_k)
\nabla^2 f_k\left(V_k^{-1/2}m_k+\widetilde Z_k\right)
\right)
\mid
\mathcal F_{k-1}
\biggr]
\\
&-
E\left[
(V_k^{-1/2}m_k)^\top\nabla f_k(\widetilde Z_k)
\mid
\mathcal F_{k-1}
\right]
\\
&-
E\biggl[
(V_k^{-1/2}m_k)^\top
\left(
\nabla f_k\left(V_k^{-1/2}m_k+\widetilde Z_k\right)
-
\nabla f_k(\widetilde Z_k)
\right)
\mid
\mathcal F_{k-1}
\biggr].
\end{align*}
The first term is zero by Gaussian integration by parts: conditional on
$\mathcal F_k$, the quantity $V_k^{-1/2}m_k$ is fixed and
$\widetilde Z_k\sim N(0,\Lambda_k)$ is independent of $\mathcal F_k$. The third term is zero because $f_k$ is $\mathcal F_{k-1}$-measurable as a
random function, $\widetilde Z_k$ is independent of $\mathcal F_k$, and
$E[m_k\mid\mathcal F_{k-1}]=0$.

Using the fundamental theorem of calculus, we obtain
\begin{align*}
& E\left[
h(S_k+T_{k+1})-h(S_{k-1}+T_k)
\mid
\mathcal F_{k-1}
\right]
\\
&=
E\biggl[
\operatorname{Tr}\left(
V_k^{-1/2}E(\Sigma_k)V_k^{-1/2}
\nabla^2 f_k\left(V_k^{-1/2}m_k+\widetilde Z_k\right)
\right)
\mid
\mathcal F_{k-1}
\biggr]
\\
&\quad
-
E\biggl[
(V_k^{-1/2}m_k)^\top
\left(
\nabla^2 f_k(\Theta V_k^{-1/2}m_k+\widetilde Z_k)
-
\nabla^2 f_k(\widetilde Z_k)
\right)
(V_k^{-1/2}m_k)
\mid
\mathcal F_{k-1}
\biggr]
\\
&\quad
-
E\biggl[
(V_k^{-1/2}m_k)^\top
\nabla^2 f_k(\widetilde Z_k)
(V_k^{-1/2}m_k)
\mid
\mathcal F_{k-1}
\biggr],
\end{align*}
where \(\Theta\sim {\rm Uniform}[0,1]\) is independent of all other random
variables.

Define
\[
    A_{k,h}(\mathcal F_{k-1})
    :=
    E\left[
    \nabla^2 f_k(\widetilde Z_k)
    \mid
    \mathcal F_{k-1}
    \right].
\]
Using
\[
    \Sigma_k=E[m_km_k^\top\mid\mathcal F_{k-1}],
\]
the last term becomes
\[
    -
    \operatorname{Tr}\left(
    V_k^{-1/2}\Sigma_k V_k^{-1/2}
    A_{k,h}(\mathcal F_{k-1})
    \right).
\]
By statement (1) of Lemma~\ref{lem: regularity}, the second term above is
bounded in absolute value by
\[
\tilde C_1(d,\beta)\|V_k^{1/2}\|_{op}
E\left[
\|V_k^{-1/2}m_k\|^{2+\beta}
\mid
\mathcal F_{k-1}
\right].
\]
Therefore,
\begin{align*}
& E\left[
h(S_k+T_{k+1})-h(S_{k-1}+T_k)
\mid
\mathcal F_{k-1}
\right]
\\
&\le
E\biggl[
\operatorname{Tr}\left(
V_k^{-1/2}E(\Sigma_k)V_k^{-1/2}
\nabla^2 f_k\left(V_k^{-1/2}m_k+\widetilde Z_k\right)
\right)
\mid
\mathcal F_{k-1}
\biggr]
\\
&\quad
-
\operatorname{Tr}\left(
V_k^{-1/2}\Sigma_k V_k^{-1/2}
A_{k,h}(\mathcal F_{k-1})
\right)
\\
&\quad
+
\tilde C_1(d,\beta)\|V_k^{1/2}\|_{op}
E\left[
\|V_k^{-1/2}m_k\|^{2+\beta}
\mid
\mathcal F_{k-1}
\right].
\end{align*}

Next, add and subtract
\[
\operatorname{Tr}\left(
V_k^{-1/2}E(\Sigma_k)V_k^{-1/2}
A_{k,h}(\mathcal F_{k-1})
\right).
\]
This gives
\begin{align*}
& E\left[
h(S_k+T_{k+1})-h(S_{k-1}+T_k)
\mid
\mathcal F_{k-1}
\right]
\\
&\le
\operatorname{Tr}\left(
A_{k,h}(\mathcal F_{k-1})
V_k^{-1/2}
\left(E(\Sigma_k)-\Sigma_k\right)
V_k^{-1/2}
\right)
\\
&\quad
+
E\biggl[
\operatorname{Tr}\left(
D_k
\left[
\nabla^2 f_k\left(V_k^{-1/2}m_k+\widetilde Z_k\right)
-
\nabla^2 f_k(\widetilde Z_k)
\right]
\right)
\mid
\mathcal F_{k-1}
\biggr]
\\
&\quad
+
\tilde C_1(d,\beta)\|V_k^{1/2}\|_{op}
E\left[
\|V_k^{-1/2}m_k\|^{2+\beta}
\mid
\mathcal F_{k-1}
\right],
\end{align*}
where
\[
    D_k:=V_k^{-1/2}E(\Sigma_k)V_k^{-1/2}.
\]

It remains to bound the weighted trace term. Since
\[
    V_k=E(\Sigma_k)+V_{k+1},
\]
we have
\[
    D_k=I-V_k^{-1/2}V_{k+1}V_k^{-1/2}.
\]
Hence
\[
    0\preceq D_k\preceq I.
\]
Therefore, for any \(x,u\in\mathbb R^d\),
\begin{align*}
&\left|
\operatorname{Tr}\left(
D_k
\left[
\nabla^2 f_k(x+u)-\nabla^2 f_k(x)
\right]
\right)
\right|
\\
&\le
\operatorname{Tr}(D_k)
\left\|
\nabla^2 f_k(x+u)-\nabla^2 f_k(x)
\right\|_{op}
\\
&\le
   \operatorname{Tr}(D_k)\,\tilde C_1(d,\beta)\|V_k^{1/2}\|_{op}\|u\|^\beta.
\end{align*}
Applying this bound with
\[
    x=\widetilde Z_k,
    \qquad
    u=V_k^{-1/2}m_k,
\]
and taking conditional expectation gives
\begin{align*}
&\left|
E\biggl[
\operatorname{Tr}\left(
D_k
\left[
\nabla^2 f_k\left(V_k^{-1/2}m_k+\widetilde Z_k\right)
-
\nabla^2 f_k(\widetilde Z_k)
\right]
\right)
\mid
\mathcal F_{k-1}
\biggr]
\right|
\\
&\le
\tilde C_1(d,\beta)\operatorname{Tr}(D_k)\|V_k^{1/2}\|_{op}
E\left[
\|V_k^{-1/2}m_k\|^\beta
\mid
\mathcal F_{k-1}
\right].
\end{align*}

Thus,
\begin{align*}
& E\left[
h(S_k+T_{k+1})-h(S_{k-1}+T_k)
\mid
\mathcal F_{k-1}
\right]
\\
&\le
\operatorname{Tr}\left(
A_{k,h}(\mathcal F_{k-1})
V_k^{-1/2}
\left(E(\Sigma_k)-\Sigma_k\right)
V_k^{-1/2}
\right)
\\
&\quad
+
\tilde C_1(d,\beta)\operatorname{Tr}(D_k)\|V_k^{1/2}\|_{op}
E\left[
\|V_k^{-1/2}m_k\|^\beta
\mid
\mathcal F_{k-1}
\right]
\\
&\quad
+
\tilde C_1(d,\beta)\|V_k^{1/2}\|_{op}
E\left[
\|V_k^{-1/2}m_k\|^{2+\beta}
\mid
\mathcal F_{k-1}
\right].
\end{align*}

Taking expectations, summing over \(k\), substituting into the Lindeberg
decomposition, and using the rescaling identity for the Wasserstein distance,
we obtain, for the fixed \(h\in Lip_1\),
\begin{align*}
& E\left[
h(W_n)-h(n^{-1/2}V_1^{1/2}Z)
\right]
\\
&\le
\frac{1}{\sqrt n}\sum_{k=1}^n
\Biggl[
\tilde C_1(d,\beta)\|V_k^{1/2}\|_{op}
E\left(\|V_k^{-1/2}m_k\|^{2+\beta}\right)
\\
&\qquad\qquad
+
\tilde C_1(d,\beta)\operatorname{Tr}(D_k)\|V_k^{1/2}\|_{op}
E\left(\|V_k^{-1/2}m_k\|^\beta\right)
\Biggr]
\\
&\quad
+
\frac{1}{\sqrt n}
\left|\sum_{k=1}^n
E\left[
\operatorname{Tr}\left(
A_{k,h}(\mathcal F_{k-1})
V_k^{-1/2}
\left(E(\Sigma_k)-\Sigma_k\right)
V_k^{-1/2}
\right)
\right]
\right|.
\end{align*}
Taking the supremum over \(h\in Lip_1\) proves the claimed bound.

Finally, statement (3) of Lemma~\ref{lem: regularity} gives
\[
    \|A_{k,h}(\mathcal F_{k-1})\|_{op}
    \le
    C\|V_k^{1/2}\|_{op}.
\]
This completes the proof.
\end{proof}

We first state the technical estimate needed to control the covariance
interaction term.

\begin{lemma}
\label{lem: H_variation}
Let
\[
    V_k:=\sum_{j=k}^n E(\Sigma_j)=(n-k+1)\Sigma_\infty,
    \qquad k=1,\ldots,n,
\]
and
\[
    H_k:=V_k^{-1/2}A_{k,h}(\mathcal F_{k-1})V_k^{-1/2},
\]
where
\[
    A_{k,h}(\mathcal F_{k-1})
    :=E[\nabla^2f_{k,h}(\widetilde Z_k)\mid\mathcal F_{k-1}]
\]
is the matrix appearing in Theorem~\ref{thm: martingale CLT}. Then, for every
$\beta\in(0,1)$,
\[
    \sup_{h\in Lip_1}\sum_{k=2}^nE\|H_k-H_{k-1}\|_{op}
    \le \frac{C}{(1-\beta)^2}n^{(1-\beta)/2},
\]
where $C$ depends on the chain, $r$, $d$, and $\Sigma_\infty$, but not on
$n$ or $\beta$.
\end{lemma}

\begin{proof}
Fix $h\in Lip_1$. Let $c_t:=\sqrt{1-e^{-2t}}$, and let
\[
    P_Ah(x):=E[h(x+A^{1/2}Z)],\qquad Z\sim N(0,I),
\]
denote Gaussian convolution with covariance $A$. The Mehler representation
of the standard normal Stein solution, followed by differentiation under the
integral, gives
\begin{equation}\label{eq: Hk_Mehler_representation}
    H_k=-\int_0^\infty e^{-2t}
    E_{\xi\sim N(0,V_{k+1})}
    [\nabla^2P_{c_t^2V_k}h(S_{k-1}+e^{-t}\xi)]\,dt.
\end{equation}
The differentiation is justified by the standard Gaussian-smoothing bounds
\[
    \|\nabla^jP_Ah\|_{op}
    \le C\|A^{-1/2}\|_{op}^{j-1},\qquad j=2,3,4.
\]
Here $\|\cdot\|_{op}$ denotes the usual matrix operator norm for matrices and Hessians, and the multilinear operator norm for high-order derivatives.


The second- and third-derivative bounds imply, respectively,
\[
\begin{split}
\|\nabla^2P_Ah(x)-\nabla^2P_Ah(y)\|_{op}
&\le 2C\|A^{-1/2}\|_{op},\\
\|\nabla^2P_Ah(x)-\nabla^2P_Ah(y)\|_{op}
&\le C\|A^{-1/2}\|_{op}^2\|x-y\|.
\end{split}
\]
Since $\min\{1,u\}\le u^\beta$ for $u\ge0$, interpolation gives
\begin{equation}\label{eq: uniform_fractional_smoothing}
\|\nabla^2P_Ah(x)-\nabla^2P_Ah(y)\|_{op}
\le C\|A^{-1/2}\|_{op}^{1+\beta}\|x-y\|^\beta,
\end{equation}
with $C$ uniform over $\beta\in(0,1)$.

Because the state space is finite, $m_k$ is uniformly bounded. Moreover,
\[
    c(n-k+1)I\preceq V_k\preceq C(n-k+1)I.
\]
Using \eqref{eq: Hk_Mehler_representation}, compare $H_k$ and $H_{k-1}$ by
changing separately the spatial shift, the covariance inside the Gaussian
convolution, and the covariance of the outer Gaussian average. Denote the
corresponding contributions by $\Delta_{{\rm shift},k}$,
$\Delta_{{\rm int},k}$, and $\Delta_{{\rm ext},k}$.

Since $S_{k-1}-S_{k-2}=m_{k-1}$, \eqref{eq: uniform_fractional_smoothing}
gives
\begin{align*}
E\|\Delta_{{\rm shift},k}\|_{op}
&\le C E\|m_{k-1}\|^\beta
\int_0^\infty e^{-2t}\|(c_t^2V_k)^{-1/2}\|_{op}^{1+\beta}dt\\
&\le \frac{C}{1-\beta}(n-k+1)^{-(1+\beta)/2},
\end{align*}
where
\[
    \int_0^\infty e^{-2t}c_t^{-(1+\beta)}dt=\frac1{1-\beta}.
\]

For the internal covariance change, stationarity gives
$V_{k-1}-V_k=\Sigma_\infty$. Along
$A_s=c_t^2(V_k+s\Sigma_\infty)$, the Gaussian-semigroup identity and the
fourth-derivative bound yield
\begin{align*}
E\|\Delta_{{\rm int},k}\|_{op}
&\le C\int_0^\infty e^{-2t}c_t^2
\|(c_t^2V_k)^{-1/2}\|_{op}^3\|\Sigma_\infty\|_{op}\,dt\\
&\le C(n-k+1)^{-3/2}\int_0^\infty e^{-2t}c_t^{-1}dt\\
&\le C(n-k+1)^{-3/2}.
\end{align*}

For the external covariance change, couple
\[
\xi_{k-1}=\xi_k+\eta_k,\qquad
\xi_k\sim N(0,V_{k+1}),\quad
\eta_k\sim N(0,\Sigma_\infty),
\]
with independent summands. Applying
\eqref{eq: uniform_fractional_smoothing} to the additional shift
$e^{-t}\eta_k$ gives
\[
E\|\Delta_{{\rm ext},k}\|_{op}
\le \frac{C}{1-\beta}(n-k+1)^{-(1+\beta)/2}.
\]
Consequently,
\[
E\|H_k-H_{k-1}\|_{op}
\le \frac{C}{1-\beta}(n-k+1)^{-(1+\beta)/2}
+C(n-k+1)^{-3/2}.
\]
Finally,
\[
\sum_{q=1}^{n-1}q^{-(1+\beta)/2}
\le \frac{C}{1-\beta}n^{(1-\beta)/2},
\]
whereas $\sum_{q\ge1}q^{-3/2}<\infty$. Summing over $k$ and taking the
supremum over $h\in Lip_1$ proves the result.
\end{proof}

\begin{proof}
The reward Poisson equation gives
\[
U_n=W_n+\frac{V(X_0)-V(X_n)}{\sqrt n},
\qquad
W_n:=\frac1{\sqrt n}\sum_{k=1}^n m_k.
\]
Since $V$ is bounded,
\[
d_{\mathcal W}(U_n,\Sigma_\infty^{1/2}Z)
\leq
d_{\mathcal W}(W_n,\Sigma_\infty^{1/2}Z)+O(n^{-1/2}).
\]
In stationarity,
\[
E(\Sigma_k)=\Sigma_\infty,
\qquad
V_k=(n-k+1)\Sigma_\infty,
\]
so the Gaussian target in Theorem~\ref{thm: martingale CLT} is exactly
$\Sigma_\infty^{1/2}Z$.

The martingale increments are uniformly bounded, and their moments can be
bounded uniformly over $\beta\in(0,1)$. Moreover,
\[
\tilde C_1(d,\beta)\leq \frac{C}{1-\beta}
\quad\text{and}\quad
\sum_{q=1}^n q^{-(1+\beta)/2}
\leq \frac{C}{1-\beta}n^{(1-\beta)/2}.
\]
It follows that the first Taylor remainder in
Theorem~\ref{thm: martingale CLT} is at most
\[
C(1-\beta)^{-2}n^{-\beta/2}.
\]
For the second Taylor remainder,
\[
D_k=\frac1{n-k+1}I,
\qquad
\operatorname{Tr}(D_k)=\frac{d}{n-k+1},
\]
and the same bound follows.

It remains to control $\Delta_n$. Let
\[
B(x):=\Sigma_\infty-v(x).
\]
Since $\pi(B)=0$, there is a bounded matrix-valued solution $\Phi$ of
\[
\Phi-P\Phi=B.
\]
We may normalize this solution so that $\pi(\Phi)=0$. Define
\[
\Gamma_k
:=
\Phi(X_k)-E[\Phi(X_k)\mid\mathcal F_{k-1}].
\]
Then
\[
E[\Gamma_k\mid\mathcal F_{k-1}]=0,
\]
and the Markov property and the Poisson equation give
\begin{align*}
B(X_{k-1})
&=
\Phi(X_{k-1})-P\Phi(X_{k-1})
\\
&=
\Phi(X_{k-1})-\Phi(X_k)+\Gamma_k.
\end{align*}
Because the chain is stationary,
\[
E(\Sigma_k)=\Sigma_\infty,
\qquad
\Sigma_k=v(X_{k-1}).
\]
Hence
\begin{equation}
\label{eq: stationary_covariance_poisson_decomposition}
E(\Sigma_k)-\Sigma_k
=
\Phi(X_{k-1})-\Phi(X_k)+\Gamma_k.
\end{equation}

For fixed $h\in Lip_1$, let
\[
H_k
:=
V_k^{-1/2}
A_{k,h}(\mathcal F_{k-1})
V_k^{-1/2}.
\]
Since $H_k$ is $\mathcal F_{k-1}$-measurable,
\[
E\operatorname{Tr}(H_k\Gamma_k)=0.
\]
Using
\eqref{eq: stationary_covariance_poisson_decomposition}
and summation by parts,
\begin{align*}
&\sum_{k=1}^n
E\operatorname{Tr}
\left(
H_k(E(\Sigma_k)-\Sigma_k)
\right)
\\
&=
E\operatorname{Tr}(H_1\Phi(X_0))
-
E\operatorname{Tr}(H_n\Phi(X_n))
\\
&\quad+
\sum_{k=2}^n
E\operatorname{Tr}
\left(
(H_k-H_{k-1})\Phi(X_{k-1})
\right).
\end{align*}
The bound on $A_{k,h}$ in
Theorem~\ref{thm: martingale CLT} and the identity
$V_k=(n-k+1)\Sigma_\infty$ imply
\[
\|H_k\|_{op}
\leq
C(n-k+1)^{-1/2}.
\]
Since $\Phi$ is bounded, the two boundary terms are $O(1)$,
uniformly over $h\in Lip_1$. Lemma~\ref{lem: H_variation} gives
\[
\sup_{h\in Lip_1}
\sum_{k=2}^n
E\|H_k-H_{k-1}\|_{op}
\leq
\frac{C}{(1-\beta)^2}n^{(1-\beta)/2}.
\]
Consequently,
\begin{align*}
\Delta_n
&\leq
\frac{C}{\sqrt n}
\left(
1+
\frac{1}{(1-\beta)^2}n^{(1-\beta)/2}
\right)
\\
&\leq
\frac{C}{(1-\beta)^2}n^{-\beta/2}.
\end{align*}

Combining the reward boundary estimate, the two Taylor-remainder
estimates, and the bound on $\Delta_n$ proves the first assertion.
Choosing
\[
\beta=1-\frac{2}{\log n}
\]
for all sufficiently large $n$ proves the second.
\end{proof}

\bibliographystyle{plainnat}
\bibliography{references}
\end{document}
